\documentclass[18pt]{article}
\usepackage{amsmath,amsfonts,amssymb}
\usepackage{natbib}
\usepackage{booktabs}
\usepackage[utf8]{inputenc}
\usepackage[T1]{fontenc}
\usepackage{amssymb}
\usepackage[version=4]{mhchem}
\usepackage{stmaryrd}
\usepackage{bbold}
\usepackage{nopageno}
\usepackage{hyperref}
\usepackage{bookmark}
\author{Adam Handwerger}
\pagestyle{plain}
\begin{document}

\noindent{\LARGE{2. \raisebox{0.4ex}{$\chi$ }= $\mathbb{R}^2$}}\\
\newline
\noindent{let x = $(x_1,x_2)$, y = $(y_1,y_2)$ and let $\phi(x)$ = $(x_1^2,\sqrt2 x_1y_1,x_2^2).$}\\
\begin{flushleft}
1. Verify that $\phi(x)^T\phi(y)$ = $(x^Ty)^2$.\\
\vspace{.5cm}
$(x_1y_1 + x_2y_2)^2 = (x_1y_1)^2 + (x_2y_2)^2 + 2x_1x_2y_1y_2$\\
\vspace{.5cm}
$\phi(x)^T\phi(y) = (x_1^2, \sqrt2 x_1x_2, x_2^2) \cdot (y_1^2, \sqrt2 y_1y_2, y_2^2)=$\\
$(x_1y_1)^2 + 2x_1x_2y_1y_2 + (x_2y_2)^2 = (x_1y_1 + x_2y_2)^2$\\
\vspace{.5cm}
2. (a) Find $\phi(x):\mathbb{R}^2 \mapsto {R}^6$ s.t. for any $(x,y)$, $ \phi(x)^T\phi(y)$=
$(x^Ty+1)^2$.\\
\vspace{.5cm}
Since $(x^Ty+1)^2$ = $x_{1}^2y_1^2+x_2^2y_2^2+2x_1x_2y_1y_2+2x_1y_1+2x_2y_2+1$, take\\
$\phi(x)$ = $(x_1^2,x_2^2,\sqrt{2}x_1x_2,\sqrt{2} x_1,\sqrt{2} x_2 ,1)$\\
\vspace{.5cm}
(b) Find $\phi(x):\mathbb{R}^2 \mapsto {R}^9$\\
\vspace{.5cm}
Take $\phi(x)$ =\\
$(x_1^2,x_2^2,\sqrt{2}x_1x_2,\sqrt{2} x_1,1\sqrt{2}x_2 ,1/ \sqrt{5},1/ \sqrt{5},1/ \sqrt{5},1/ \sqrt{5},1/ \sqrt{5})$.
\vspace{.5cm}
Verify that $K_{l}(x_i, x_j)=\left(1+x^{\top} y\right)^{d}$ is pos. definte for $d=1,2, \ldots$\\
\vspace{.5cm}
(a) $K_{d}(x_i,x_j)$ is symmetric since $K_{d}(x_i,x_j)$ = $K_{d}(x_j, x_i)$ by the sym property of the inner product\\
\vspace{.5cm}
(b) $K_{1}(x, y) :=K$ is positive definite since for any $v\in \mathbb{R}^{n}$ we have

$$
\begin{aligned}
&v^{\top} K v=\sum_{i j}^n v_{i} v_{j} k_(x_i,x_j)=\sum_{i j} v_{i} v_{j}\left(1+x_{i}^{\top} x_{j}\right) \\
= & \sum_{i} \sum_{j} v_{i} v_{j}+\sum_{i} \sum_{j} v_{i} v_{j} x_{i}^{\top} x_{j} \\
= & \sum_{i} v_{i} \sum_{j} v_{j}+\left(\sum_{i} v_{i} x_{i}^{\top}\right)\left(\sum_{j} v_{j} x_{i}\right) \\
= & \left(\sum_{k} v_{k}\right)^{2}+\left(\sum_{i} v_{i} x_{i}\right)^{\top}\left(\sum_{j} v_{j} x_{i}\right)  \geqslant 0\\
\end{aligned}
$$
If $w:=\sum_{i} v_{i} x_{i}$ then the sum on the right is the norm of w and so both terms are greater or equal to zero. 
\pagebreak
(b) continued

Assume \(k_{d-1}(x, y)=\left(1+x^{\top} y\right)^{d-1}\) is pos. definite.\\
Then \(k_{d}(x, y)=\left(1+x^{\top} y\right)^{d}=\left(1+x^{T} y\right)^{d-1}\left(1+x^{\top} y\right)\) is also pos. definite
by (c) below. Since \(\left(1+x^{\top} y\right)\) is pos definite, by induction \(K_{d}(x, y)=\left(1+x^{\top} y\right)^{d}\)
is also.\\
\vspace{.5cm}
(c) If \(k_{1}(x, y)\) and \(k_{2}(x, y)\) are pos. definite that so is their product \(k=k_{1} \cdot k_{2}\)
\vspace{.5cm}
Proof.\\
\vspace{.5cm}
Define the following matrices:\\
\(C\) is the matrix with elements, \(c_{i j}=k_{1}\left(x_{i}, x_{j}\right)\),\\
\(D\) is the matrix with elements, \(d_{i j}=k_{2}\left(x_{i}, x_{2}\right)\)\\
\(E\) as the matrix formed by the elements, \(e_{i j}=c_{i j} d_{i j}\),\\
\vspace{.5cm}
E is clearly symmetric since \(k_{1}\) and \(k\), are both symmetric\\
\vspace{.5cm}
Pick and \(u \in \mathbb{R}^{2}\). Since \(C\) and \(D\) are pos. definite they can be decomposed as

\[
C=A^{\top} A \text { and } D=B^{\top} B
\]

for some matrices \(A, B\).\\
\vspace{.5cm}
Now write

\[
A=\left(\begin{array}{llll}
a_{1} & a_{2} & \cdots & a_{n}
\end{array}\right), B=\left(\begin{array}{llll}
    b_{1} & b_{2} & \cdots & a_{n}
    \end{array}\right)
\]

where \(a_{i}\) and \(b_{i}\) are the columns of \(A, B\) respectively. Therefore,\\

\[
c_{i j}=a_{i}^{\top} a_{j}=\sum_{k} a_{i k} a_{j k}, \quad d_{i j}=b_{i}^{\top} b_{j}=\sum_{l} b_{i l} b_{j l}
\]

Finally,

\[
\begin{aligned}
u^{\top} E u & =\sum_{i j} u_{i} u_{j} c_{i j} d_{i j}=\sum_{i j} u_{i} u_{j}\left(\sum_{k} a_{i k} a_{j k}\right)\left(\sum_{l} b_{i l} b_{j l}\right) \\
& =\sum_{k l} \sum_{i j} u_{i} u_{j} a_{i k} a_{j k} b_{i l} b_{j l} \\
& =\sum_{k l}\left(\sum_{i} u_{i} a_{i k} b_{i l}\right)\left(\sum_{j} u_{j} a_{j k} b_{j l}\right)=\sum_{k l} f_{k l}^{2} \geqslant 0
\end{aligned}
\]

(since Dummy indices don't care what you call them)
\pagebreak

(d) Let $x=(x_1,x_2)$ and\\
    \begin{align}
    \phi(x)= 
    \left\{
        \begin{aligned}
             & 1, d=0\\
             & (\phi_{d-1}(x),x_1 \phi_{d-1}(x),x_2 \phi_{d-1}(x))^{\top}, d=1,\dots\\             
        \end{aligned}
    \right.    
    \end{align}

    


For $d = 1$, the kernel $K_d(x,y)$ = $\phi(x)_d^{\top}\phi(y)_d$ = $(1+x^{\top}y)$ is well defined,\\
see (a) above. For $d > 1$ the feature map can be written as $\phi_d(x) = (1, x_1,x_2) \phi_{d-1}(x)$.
Assume that $\phi(x)_d^{\top}\phi(y)_d$ is well defined. Then\\
since $\phi(x)_{d+1}^{\top}\phi(y)_{d+1}$$\left(1+x^{\top}y\right) \phi(x)_d^{\top}\phi(y)_d$ by (c) above and induction\\
$K_d(x,y)$ is well-defined for all d.\\
\vspace{.5cm}
4. (a) Let $K(x, y)$ be pos. definite. Then for any u $\in \mathbb{R}^{2}$\\
$u^{\top} k(x, y) u \geqslant 0\Rightarrow \text { for } k^{\prime}(x, y)=-k(x, y)\Rightarrow u^{\top} k^{\prime}(x, y) u \leq 0$\\
\vspace{.5cm}
(b) Find a $\phi(x): \mathbb{R}^{2} \rightarrow H \text{ where for any } x,y $

\[
\phi(x)^{\top}\phi(x)=x^{\top} y-1 .
\]

Proof: None exist Since this must hold for \(x=(0,0)\) and \(y=(0,0)\)\\
which would imply that\\
\newcommand{\norm}[1]{\left \lVert #1 \right \rVert}
\hspace{4.2cm} $\phi(x)^{\top} \phi(x)$ = $\norm{\phi(x)}^2=-1$\\
\vspace{.5cm}
This is impossible since $\|\cdot\| \geqslant 0$
\end{flushleft}
\end{document}